
Cross-model benchmark correlations are a common tool for characterizing alignment structure: when TruthfulQA and BBQ co-move positively across a population of open-weight LLMs, researchers may infer that the two constructs --- factuality and bias avoidance --- reflect a shared alignment property. This inference is potentially confounded by general model capability. Larger, more capable models tend to score higher on all benchmarks simultaneously, producing positive cross-model correlations that may reflect scale rather than any alignment-specific training signal.

The raw Spearman correlation between TruthfulQA MC2 and our bias proxy across N=296 open-weight LLMs is 0.732. After controlling for MMLU via partial Spearman, this correlation decreases to 0.343 --- a 53.1\% reduction. The Fisher z test between raw and partial estimates yields z = 6.97, p = 3.22 $\times 10^{-12}$. More than half of the apparent alignment co-movement is attributable to MMLU-indexed general capability.

This finding has methodological consequences. Alignment benchmark evaluation in the open LLM evaluation literature --- exemplified by leaderboard-based comparisons and technical reports such as Llama 2 \citep{Touvron2023} --- typically reports raw pairwise correlations or paired improvements without controlling for scale. Positive cross-model correlations between factuality, bias avoidance, and safety benchmarks are sometimes interpreted as evidence that alignment training generalizes across dimensions. Our results indicate that a substantial fraction of such co-movement is explained by MMLU, not by any alignment-specific mechanism.

Prior work has addressed the global dimensionality of benchmark suites. BenchScope \citep{Sha2026} reports an effective dimensionality of approximately 1.7 for the Open LLM Leaderboard, consistent with a small number of latent evaluation axes. PCA-based analyses \citep{clawrxiv2603} show that TruthfulQA loads primarily on a second principal component orthogonal to general capability (PC1). Neither line of work provides a directional, pairwise hypothesis test for whether a specific scale covariate (MMLU) confounds a specific benchmark pair (TruthfulQA $\times$ BBQ), nor does it quantify the magnitude of that confound via a formal statistical test.

We address this gap by applying the Fisher z difference test between raw and partial Spearman correlation as a statistical diagnostic. The partial Spearman correlation controlling for MMLU directly measures alignment-specific co-movement after removing scale-driven variance; the Fisher z test provides a formal check of whether MMLU control changed the correlation significantly.

The contributions of this paper are as follows:

\textbf{(1)} We quantify scale confound magnitude: MMLU explains 49\% of TruthfulQA variance and 76\% of bias-proxy variance across N=296 open-weight LLMs, and MMLU control reduces the pairwise correlation by 53.1\% (Fisher z p = 3.22 $\times 10^{-12}$, BCa CIs non-overlapping).

\textbf{(2)} We demonstrate the Fisher z raw-vs-partial difference test as a practical diagnostic for detecting capability-scale confounding in alignment benchmark correlation analysis.

\textbf{(3)} We report a positive residual partial correlation (0.343, p = 1.40 $\times 10^{-9})$ after scale removal, indicating non-zero alignment-specific co-movement, while noting that definitive structural interpretation is not possible at N=296 given a confidence interval that spans the pre-specified coherence threshold.

\textbf{(4)} We report a null RLHF sign test result on the bias proxy (p = 0.686), contrasting with established RLHF improvement on TruthfulQA of +3.406 points (BCa CI [2.589, 4.212]), and note this as a hypothesis-generating finding subject to proxy limitations.

